\setcounter{section}{0}
\setlength{\parindent}{0pt}  % niente rientro
\setlength{\parskip}{6pt}
\setlength{\itemsep}{3pt}
\section{Organizzazione del lavoro}
L'organizzazione del lavoro svolto sul progetto Rinova è stata strutturata tenendo conto della dimensione ridotta del gruppo.
Per compensare il maggior carico di lavoro i compiti sono stati suddivisi in base alle competenze tecniche
ed all'esperienza pregressa dei membri, ottimizzando le tempistiche e permettendo di portare avanti in parallelo (in modalità asincrona) sia lo sviluppo dell'applicazione web che la stesura della documentazione.

Nonostante l'autonomia operativa la coesione del team è rimasta centrale durante tutto il processo produttivo grazie ad un piano di comunicazione strutturato come segue:

\begin{itemize}
    \item Micro aggiornamenti e Revisione: Il completamento di componenti minori veniva prontamente notificato solitamente via messaggio, e veniva sottoposto quanto prima ad una rapida revisione da parte dell'altro membro del gruppo, così da individuare più facilemnte criticità oltre che a garantire che nulla venisse integrato senza il consenso comune.
    \item Checkpoint: Al completamento di macro componenti o funzionalità principali sono state effettuate delle chiamate di allineamento (solitamente attraverso l'applicazione Discord). Queste chiamate avevano lo scopo di verificare lo stato del processo, risolvere dubbi critici e se necessario ridistribuire i ruoli per le fasi successive, ridimensionando il carico di lavoro rispetto agli impegni personali dei membri del gruppo.
\end{itemize}

Tale approccio ha permesso di sopperire efficacemente alla mancanza di incontri in presenza, limitati dalla mancata divergenza di orari personali ed accademici dei membri.

Per quanto concerne gli strumenti utilizzati, il coordinamento è avvenuto tramite Discord, sfruttando le chiamate vocali per i meeting e la chat di progetto per la messaggistica istantanea La gestione del lavoro è stata centralizzata su GitHub, utilizzato per il versionamento sia del codice della WebApp che dei deliverable (redatti in LaTeX).
L'intero processo di scrittura è stato supportato dall'uso di Visual Studio code come ambiente di lavoro.
